% The integrated commentary of the concise study, beginning on physical page
% 5. It is built on questions that more than one of the three interpretations
% answers: to whom the commands of flight are spoken; what the body and the
% eagles are; what captivity the Introit brings back; why deliverance is sung
% in the past tense, with the Epistle's two tenses; what the fig tree, this
% generation and the words that do not pass assure; whose cry the chants of
% the depths are; what the Communion promises and to whom; and what the three
% orations ask. Under each, the authorities of the different interpretations
% are set side by side at the text they share. A short comparison of the four
% senses closes it. Every claim, attribution and qualification is the
% expansive study's; the connecting prose is written for this companion.

\runninghead{Detailed Commentary}
\section{The Propers: Detailed Commentary}
\zlabel{triptych:concise:commentary:start}

\subsection{To whom the commands of flight are spoken}

Chrysostom holds to the letter and gives his reasons. The sabbath shows whom
the Lord is addressing: ``Seest thou that His discourse is addressed to the
Jews \dots\ the apostles surely were not to keep the Sabbath day, neither to
be there, when Vespasian did those things'' (\work{Hom.\ in Mt.}~76.1). The
tribulation without equal is the siege, and whoever doubts it is to ``study
the writings of Josephus.'' ``No flesh'' means ``no Jewish flesh,'' and the
elect for whose sake the days were shortened are ``The believers that were
shut up in the midst of them'' (76.2). St Thomas Aquinas follows the same
line: the believers left the city for the realm of Agrippa, and the man on
the roof and the man in the field are told to leave everything, since a man
will give all he has for his life (\work{Super Matthaeum}, cap.~24). St Bede,
on Mark's account of the discourse, records that the Christians of the
province, warned by an oracle, withdrew across the Jordan to the city of
Pella, ``as the Church History relates'' (\work{In Marci evangelium
expositio}~IV, on 13:14).

Hilary stops at the housetop. ``This precept, by human understanding, does
not admit a reason for what is said and done. For whoever is on the roof
cannot leave unless he comes down from it''; nor is anything gained by
standing on the roof rather than staying in the house (\work{In
Matthaeum}~25.5). He therefore begins from what a roof is, ``the summit of the
house, and the lofty completion of the whole dwelling.'' The one who stands
there is the one ``new by rebirth, lofty in spirit,'' and he ``ought not to go
down into lower things by desire for worldly things''; the man in the field,
at work on the commandment, is not to take up again the tunic ``of old sins.''
Since nothing is evil to a man ``except by his own fault,'' the woe to women
with child cannot be a curse on a condition of nature. The pregnant are souls
heavy with sins, and the nursing are souls still fed on milk, who ``have
neither escaped their sins nor received the food of the true bread'' (25.6).
Winter and sabbath mean ``that we be not found either in the cold of sins or
in idleness from good works'' (25.7). The abomination itself Hilary takes for
Antichrist, and the flight from Judea for a flight from the contagion of those
who will receive him (25.3--4).

Jerome reaches a like rule from the abomination. After its plain senses he
adds that it ``can be understood also as every perverse doctrine.'' When such
doctrine stands in the holy place, ``that is, in the Church,'' we must flee
to the mountains, ``leaving the letter that kills, to draw near to the
eternal mountains'' (Ps 75:5), stay on the roof, where the devil's fiery darts
cannot reach, and not go down ``from the house of our former life'' (on
24:15--18). He reads the winter and the sabbath first of the siege under
Titus and Vespasian, when cold and the sabbath law would hinder flight, and
then of the end, as a command ``that our faith and our charity towards Christ
not grow cold, and that we not lie idle and sluggish in the work of God, on a
sabbath from the virtues'' (on 24:20). Bede, after the history of Pella,
gives a spiritual sense of his own that runs close to Jerome's. When heresies
and crimes reign among those who seemed consecrated to the heavenly
mysteries, those who stand in the confession of the true faith ``ought to
climb so much the higher to the summit of the virtues, as we see more
following the broad and wandering ways of the vices.'' The man on the
housetop is one who, ``rising in mind above carnal things, lives spiritually
as in free air,'' and is not to ``seek again the desires of the world or the
flesh which he had left''; the house he would go down into is, for Bede,
either this world or the flesh (on 13:15--16). Augustine gives the winter and
the sabbath a second sense of the same kind: ``that no one be found by that
day in sorrow or joy over temporal things'' (\work{Quaestiones
Evangeliorum}~I.37).

The difficulty is one of method, and the Fathers do not settle it. Hilary
leaves the letter of the housetop because, he argues, it has no reason;
Chrysostom does not leave the letter, since the words were spoken to Jews who
would see the city taken, and the apostles would not be there. On
Chrysostom's account the passage is a fulfilled history, read on this Sunday
as proof that what the Lord foretold came true. On Hilary's and Jerome's it is
a command to the congregation itself, who are on the roof by baptism and in
the field at their work. The second interpretation follows Hilary and Jerome
and sets Chrysostom's sense first, as Jerome sets out the plain senses before
the spiritual one and Bede tells of Pella before the climb to the summit of
the virtues: the history shows that the command was obeyed and that those who
obeyed it were saved, and the spiritual sense applies the same command to
those who hear it now. Hilary and Jerome still differ on the abomination:
Hilary names Antichrist, while Jerome leaves several referents open and adds
every perverse doctrine.

\subsection{The body and the eagles}

The proverb of verse~28 is about carrion birds and a carcass, and it stands
between the lightning of the coming and the darkening of the sun. Jerome,
Hilary, Augustine and Chrysostom agree that the body is Christ's and the
eagles are the saints, and each gives a different reason for reading it so.

Jerome's reason is the word. Eagles sense a carcass even across the sea; how
much more, he argues, must we ``hasten to him whose lightning goes out from
the east.'' The Greek has \latin{ptoma}, a fallen thing, ``which in Latin is
more expressively called \latin{cadaver} because it falls by death,'' and so
the body can be understood ``as the Passion of Christ, to which we are
summoned, so that wherever it is read in the Scriptures we may gather, and
through it may come to the Word of God'' (on 24:28). His eagles are ``the
saints, whose youth is renewed like the eagle's'' (Ps 102:5), who take wings
(Is 40:31) ``so as to come to the Passion of Christ.'' Hilary's reason is the
fitness of the place. The Lord called the saints eagles ``from the flight of
the spiritual body,'' and showed that their assembly, when the angels gather
them, ``will be in the place of the Passion. And it is fitting that the coming
of his brightness be awaited there, where he worked for us the glory of
eternity by the passion of his bodily lowliness'' (25.8). Augustine's reason
is the speaker. The body is the one the Lord raised with him into heaven, ``and
it was called a carcass because he was about to die when he said these
things''; the eagles are ``the spiritual, who, by imitating his Passion and
lowliness, are fed as it were from his body'' (\work{Quaestiones
Evangeliorum}~I.42). His ``as it were'' keeps the feeding a likeness: they are
nourished by imitating the Passion. Chrysostom's reason is the context. The
coming ``needs not a herald'' but shows itself everywhere at once, like
lightning, and the eagles are ``the multitude of the angels, of the martyrs,
of all the saints''; the sign of the Son of man is ``the cross being brighter
than the sun'' (\work{Hom.\ in Mt.}~76.3).

They do not agree on the place: the Passion wherever it is read, the place of
the Passion, heaven, wherever Christ appears. None of the four excludes the
others. Jerome, Hilary and Aquinas name the Passion as what draws the
gathering; Aquinas says that the word \latin{cadaver} signifies it ``because
Christ will then come showing the signs of his Passion'' (\work{Super
Matthaeum}, cap.~24). Hilary reads the signs in the sky as showing ``the glory
of his coming and the return of his brightness,'' and after likenesses of
lowly things made great, the great tree from the mustard seed among them, he
ends, ``so, from the lowliness of death, the glory of God''; by them the Lord willed
us to know ``the time of our redemption, in which we shall be carried over
from the corruption of bodies into the honour of a spiritual substance''
(26.1). His verb for that passage, \latin{transferemur}, is the verb the
Epistle uses in the past tense when it thanks the Father who
\latin{tránstulit}, ``hath translated'' us into the kingdom of his Son (Col
1:13). If their answers stand, the verse promises that the end will be an
assembly of the scattered to the one who died, and that is the first
interpretation's claim.

\subsection{Which captivity the Introit brings back}

Between ``I will bring back your captivity'' and ``from all places'' Jeremias
has ``and I will gather you out of all nations'' (Jer 29:14), and the antiphon
does not sing it. The verb it passes over is the Gospel's for the eagles,
\latin{congregabúntur}, and for the angels who \latin{congregábunt eléctos
eius}, ``shall gather together his elect'' (Mt 24:28, 31). Because God's
promise to the exiles and the Lord's account of the end use one word for the
gathering of God's people, the first interpretation takes the return from all
places as a first instance of the gathering from the four winds. Jerome,
commenting on the letter itself, gives it a sense beyond the return from
Babylon: ``According to the anagogy, we are tossed in the confusion of this
world until we are found worthy to receive the rest of the number seven, and,
our penance accepted, God fulfils what he promised, and we are brought back
to our own place, the Church'' (\work{In Hieremiam}~V, on 29:10--14). The
exiles of this sense are sinners, and the place to which they return after
penance is the Church. Augustine, on the psalm verse, asks which captivity is
meant, since ``I think that none of us is now among the barbarians.'' He
answers with Paul's ``law of sin'' (Rom 7:23): Christ turned that captivity
away ``not by freeing us from barbarians, but by freeing us from evil works,
from our sins, through which Satan ruled over us''; ``iniquity held you
captive; when iniquity is forgiven you are set free'' (\work{Enarrationes in
Psalmos}~84.4). He draws a duty from it at once: ``Confess, then, that you
are in captivity, that you may be worthy to be set free'' (84.5). Augustine
does not comment on Jeremias, and the first interpretation reads the
antiphon's captivity as he reads the verse's because the two are sung in one
chant.

Rupert of Deutz, a twelfth-century abbot who wrote on the liturgy of the
Church's year, reads the chants as the story of Israel's remnant at the end.
In the Introit God promises a banquet of grace to his brethren according to
the flesh, as Joseph, ruling Egypt, received and feasted the brothers who had
sold him (Gen 43:16, 34). The remnant will come back to the one ``whom they
denied and killed,'' and ``will be received by him into grace''
(\work{De divinis officiis}~XII.23).

The second interpretation uses another part of the same letter. Before the
promise Jeremias warned the exiles, ``Let not your prophets that are in the
midst of you, and your diviners deceive you'' (Jer 29:8), and the Gospel
warns against those who say that Christ is ``in the desert'' or ``in the
closets'' (Mt 24:26). Jerome reads the desert as the doctrine of the
philosophers and the inner rooms as the hiding places of heretics (on
24:25--26). Augustine says that the Lord ``fittingly named the Church
lightning here,'' spread through the whole world, makes the ``here'' and
``there'' the schisms that hold some part of the earth, and the desert and
the inner rooms the ``dark and hidden meetings'' of heretics
(\work{Quaestiones Evangeliorum}~I.38). The one who has gone up to the
mountains is not to go out after any of them.

The third interpretation takes the antiphon's first words, ``I think thoughts
of peace.'' Jerome sets God's knowledge against the exiles' prophets: ``He
says that he knows what he thinks, and that they, with their prophets, do not
know. The knowledge of things to come therefore belongs to God alone.'' On the
promise to hear he asks why the exiles must pray at all: ``The Lord could
certainly have fulfilled what he had promised without the calling upon him
and the prayer of the captives; but he urges them to prayer, so that they may
be worthy to receive what is promised,'' and he cites ``Ask, and it shall be
given you'' (Mt 7:7; on 29:10--14). In this interpretation the Introit is a
promise whose keeping God alone foresees, and for which he still bids his
people pray.

\subsection{Deliverance in the past tense}

The Gradual's \latin{Liberásti nos}, ``Thou hast delivered us,'' is read three
ways. For Rupert those ``freed from the captivity of souls'' give thanks in it
(XII.23). For Aquinas the deliverance is partly present already: on
\latin{salvasti \dots\ ex affligentibus nos} he writes, ``either in the
present, because though they may harm us in body, yet in spirit they cannot
harm us,'' or in the future (\work{Super Psalmos}~43, nn.~3--4). The devil's
fiery darts that Jerome keeps from the one on the roof are aimed at the
spirit, and so, in the second interpretation, the one who stays on Jerome's
roof has Aquinas's present deliverance already.

For Augustine the past tense is the form of a certain future. On the
Gradual's own verse he writes: ``This too is said of the future in the form
of the past; but it is said as past because it is as certain as if it had
been done'' (\work{Enarrationes in Psalmos}~43.8). The prophets, he explains,
foretold the Passion in past tenses, ``They have dug my hands and feet'' (Ps
21:17). ``For to us the things that are past are certain, and the things to
come uncertain''; but ``Give me a prophet to whom the future is as certain as
the past is to you'': as what you remember cannot not have happened, so for
him what he knows will be cannot not be. In the next verse, which the Gradual
also sings, the psalm changes tense, ``we shall glory \dots\ we shall give
praise,'' and Augustine notes it: ``See how he mixes in words of the future
tense too'' (43.9). Of the Introit's verse he says that the prophet ``uses
words as if of time already past, because with God even what is to come is
already done'' (84.3). The third interpretation reads the past tenses of both
chants by these sentences.

Gregory and Chrysostom give another ground. Gregory, preaching on Luke's form
of the discourse, can see many of the foretold evils in his own time, and
``since many things foretold are already fulfilled, there is no doubt that the
few that remain will follow, because the showing of past things is the
certainty of those that follow'' (\work{Homiliae in Evangelia}~I.1.1).
Chrysostom bids whoever doubts ``test His sayings, and when he hath found
them true \dots\ from what is past, let him believe also the things to come''
(\work{Hom.\ in Mt.}~77.1). They argue from what has already happened to what
will happen; Augustine argues from God, for whom the future is as fixed as the
past. The third interpretation uses both arguments.

\subsection{The Epistle: what is done and what is asked}

Each interpretation takes a different part of Paul's sentence. The first
takes the clauses of thanks. Chrysostom notices both verbs: Paul said
``delivered,'' and not ``sent forth,'' ``showing our great misery''; and he
said ``translated,'' as one leads a soldier over from one position to
another, ``so that it is both of us and of Him.'' Those ``that were enemies,
those that were in darkness'' were brought ``to where the Son is, to the same
honor with Him'' (\work{Hom.\ in Col.}~2). Aquinas glosses the power of
darkness as ``whether of demons or of sins,'' and quotes Isaias, ``even the
captivity shall be taken away from the strong'' (Is 49:25; \work{Super
Colossenses}, cap.~1, lect.~3).

The second takes the petition for the walk. Chrysostom notes that ``with
faith he always couples conduct''; that ``To walk'' means ``always, not once,
but continually''; and that bearing fruit means being ``always advancing, and
nowhere standing still.'' The strength Paul asks for concerns ``trials and
persecutions,'' and the two virtues are distinct: ``long-suffering towards
one another, patience towards those without.'' Aquinas distinguishes them
otherwise: patience is against the hardness of adversity, longsuffering
against the delay of the reward, and ``with joy'' against bearing them in
sadness; and ``for virtue it is not enough to know or to will, unless one
works without wavering.'' Fromage, on this lesson, writes that ``The just man
cannot possibly remain stationary in this world; he must either descend or
ascend'' (vol.~11, p.~505).

The third takes the sentence whole: its petitions look forward and its thanks
look back (Col 1:9--11, 12--14). Chrysostom says why Paul, about to exhort
them, ``makes no mention of what by and by shall be laid up for them'' but ``mentions the
things which were already theirs, for these are the causes of the other'':
``For that which hath already come to pass gains belief.'' Aquinas reads the
translation into the kingdom in two ways: ``And this happens when we are
freed from sin \dots\ or, according to the letter, that we might attain
eternal life''; and for longsuffering against delay he cites Habacuc: ``if
it make any delay, wait for it: for it shall surely come'' (Hab 2:3).

\subsection{The fig tree, this generation, and the words that do not pass}

The Fathers do not read the fig tree alike. For Chrysostom the parable shows
the interval to be short, and foretells ``a spiritual summer, and a calm that
should be on that day (after the present tempest) for the righteous; but to
the sinners the contrary, winter after summer.'' Above all it is given ``to
confirm His saying, as assuredly thus to come to pass. For as this of the fig
tree is of necessity, so that too,'' as the grain of wheat that must die
proves the resurrection (Jn 12:24; 77.1). Gregory, on Luke's form, where the
trees ``shoot forth their fruit'' and ``the kingdom of God is at hand'' (Lk
21:29--31), reads it by contrast: ``as from the fruit of the trees the
nearness of summer is known, so from the ruin of the world the nearness of
the kingdom of God is known.'' The world's fruit is its fall, ``For it grows
in order to fall,'' and the kingdom is summer ``because then the clouds of our
grief pass away, and the days of life shine with the brightness of the eternal
sun'' (I.1.3). Jerome reads the parable as a warning against haste: ``do not
think that the consummation of the world is already here, but that certain
preludes and forerunners, as it were, are coming'' (on 24:32--33). Augustine
makes the fig tree the human race, whose tender branch is the time when the
sons of men ``have advanced through the faith of Christ to spiritual fruits''
(\work{Quaestiones Evangeliorum}~I.39). Hilary goes the other way: the branch
is Antichrist, its leafing the green of sins in their triumph, and the summer
the heat of eternal fire (26.2). The third interpretation follows Chrysostom
and Gregory, for whom the fig tree is a pledge of summer.

On ``this generation'' the answers differ again. Chrysostom will not take it of
those then living: the Lord speaks ``not of the generation then living, but of
that of the believers,'' as the psalm says ``This is the generation of them
that seek him'' (Ps 23:6), and ``over this generation shall nothing prevail,
not famine, not pestilence'' (77.1). Jerome gives two other answers: it is
``either the whole human race, or the Jews in particular'' (on 24:34). Hilary
speaks of ``our generation,'' which cannot pass before all these things have
run their course, and gives the saying's purpose: ``so that faith in things to
come might be certain, he added, by saying Amen as a profession of truth,''
that this generation would not pass (26.3). Taken of those alive when it was
spoken, the saying would be a prediction that failed, since they did not see
the Son of man come on the clouds. None of the three takes it so, and each
leaves the day unknown, as verse~36 says it is.

On heaven and earth Gregory finds a paradox: ``nothing in the nature of
bodily things is more lasting than heaven and earth, and nothing in nature
passes so quickly as speech.'' The Lord's meaning, he says, is that whatever
seems lasting to us does not last into eternity without change, ``because my
speech, which passes, utters sentences that remain without change'' (I.1.4).
Chrysostom says that heaven and earth are named because ``it were more easy
for these firm, fixed, and immoveable bodies to be blotted out, than for ought
of my words to fall to the ground,'' and to show the speaker ``maker of all''
(77.1). Hilary finds the reason in the origin of each: heaven and earth,
``come forth from nothing,'' ``have in themselves the necessity of not
being,'' while his words, ``drawn from eternity,'' ``contain in themselves the
power to remain'' (26.3). Jerome says that heaven and earth will pass ``by
change, not by their own abolition'' (on 24:35), and Bede, on Mark's form,
that they pass in the form they now have and subsist without end in their
being (on 13:31). Hilary stresses that created things can cease, Jerome and
Bede that they will be changed; they differ in emphasis and not in doctrine.

\subsection{The cry from the depths}

The Alleluia and the Offertory sing ``From the depths I have cried to Thee, O
Lord.'' Augustine, preaching on the psalm, recalls that Jonas cried from the
whale's belly, and says: ``we also ought to understand from what deep we cry
unto the Lord. For this mortal life is our deep''; and ``they are very deep in
the deep, who do not even cry from the deep'' (\work{Enarrationes in
Psalmos}~129.1, in the abridged English of the Nicene and Post-Nicene
Fathers). On the psalm's fourth verse, which is not sung, he says that the
``pouring forth of innocent blood blotted out all the sins of the guilty: so
great a price paid down redeemed all captives from the hand of the enemy who
captured them'' (129.3). In the first interpretation, since for Augustine this
mortal life is the deep, those who cry from it are those whom the angels of
verse~31 have not yet gathered.

Hilary asks the reader to picture the psalmist ``praying from the deep surge
of human calamities,'' and ``not with the lips, but in the secret of the
heart,'' as the Lord commanded (Mt 6:6); the cry fits one who ``either does
not know, or suffers, or prays'' (\work{Tractatus super Psalmos}, on Ps 129).
The Gospel's only command to pray is ``pray that your flight be not in the
winter or on the sabbath,'' and Aquinas says that the Lord bids them pray
``because in such things one must have recourse to God alone'' (\work{Super
Matthaeum}, cap.~24). The second interpretation hears the Alleluia and the Offertory as prayer of that kind.
The third uses another sentence of Hilary's on the same psalm. God, who is
everywhere, hears in any case, ``but he is to be prayed to by us'': ``It
belongs to his nature to hear; it belongs to our faith to pray that we be
heard.''

Two medieval commentators on the liturgy give the chants a sequence. For
Rupert the Offertory is the brothers' plea to Joseph, ``I beseech thee to
forget the wickedness of thy brethren'' (Gen 50:17), and the Communion is
Joseph's answer, ``Fear not \dots\ You thought evil against me: but God turned
it into good'' (Gen 50:19--20). Sicard of Cremona, a medieval bishop who wrote
on the offices of the year, says that those freed from captivity give thanks
in the Gradual, pray in the Offertory, and ``receive in the Communion,
\emph{Amen dico vobis}, the assurance of their petition'' (\work{Mitrale}~VIII.25);
in another chapter, that ``the Lord in the Communion assures them of
salvation'' (VIII.23). Rupert had already taken the Gradual for the thanks of
those freed and the Offertory for their supplication; the Communion as
assurance is Sicard's own.

\subsection{Whatsoever you ask}

The second and third interpretations read the Communion in different ways.
Mark gives its promise in another setting, after the withered fig tree, but
its ``whatsoever'' excepts no prayer, and so, in the second interpretation,
it reaches the prayer that the Gospel commands of those
who flee. The third looks at the verb, \latin{accipiétis}, ``you shall
receive,'' which is future. Augustine's prophet holds what is to come as
certain as what is past; the third interpretation applies this to the one who
prays, who, believing that he shall receive, holds the thing asked for with
that certainty before it is given.

Read alone, the antiphon promises anything whatever to anyone who believes.
In Mark the next verse commands forgiveness, and Bede, on that verse, puts the
condition in the one who prays: ``The difference among those who pray is to
be noted.'' Whoever has perfect faith, ``which works through love,'' can move
spiritual mountains by prayer or even by command. Those not yet perfect should
ask that their sins be forgiven, ``and without doubt they will obtain what
they ask, provided that they first forgive others who sin against them''
(\work{In Marci evangelium expositio}~III, on 11:25). Three sayings on other
texts also look to the one who prays rather than to the thing asked: Hilary's
on the psalm of the Alleluia and Jerome's on Jeremias's letter, both given
above, and Aquinas's on Paul's ``we cease not to pray,'' where he defines
prayer as ``the ascent of the mind to God,'' made with devotion ``not
to bend him, but to raise ourselves to him'' (\work{Super Colossenses},
cap.~1, lect.~3). None of the three comments on Mark. Their agreement with
Bede is the third interpretation's reason for reading the Communion as a
lesson in how to pray, and not as a guarantee of whatever object is asked.

\subsection{Wills, hearts and the mind: the three prayers}

No Father cited here comments on the three orations, and each interpretation
brings to them a Father's words on another text. In the first, the haste
toward Christ that Jerome draws from the eagles starts in the will, and the
Collect asks God to stir the wills of his faithful. In the second, the Collect
asks for the opposite of what Jerome feared, that we lie idle
``in the work of God'': it uses the same noun, \latin{divíni óperis fructum},
``the fruit of the divine work,'' and asks for wills roused and that fruit
followed \latin{propénsius}, ``more readily,'' where Lasance has ``more
earnestly.'' In the third, its \latin{remédia maióra}, ``greater remedies,''
ask for more than has yet been given.

The Secret asks that God turn \latin{ómnium nostrum \dots\ corda}, ``the hearts
of all of us,'' \latin{ad te}, ``to thee.'' The first interpretation reads that
as the hearts of many turned to one God, as the elect of verse~31 are brought
from every quarter to one Lord. The second notes that Hilary's man on the
housetop must not go down by \latin{cupiditate}, ``desire,'' for worldly
things, and that Bede's must not seek again the \latin{desideria} of the
world and the flesh that he has left. Both commands assume that the one who
keeps them has already left his former desires, and the Secret asks that
leaving for all of us, with the same two nouns: delivered from earthly
\latin{cupiditátibus}, we are to pass over to heavenly \latin{desidéria}.
Gregory, on the same discourse in Luke, asks, ``What is mortal life but a
way?'', and wonders at those who grow weary on the way and yet do not wish it
to end (I.1.3); at the end of the sermon he counsels, ``Do not then, my
brethren, love what you see cannot stand for long'' (I.1.5, citing 1~Jn 2:15).
In the third, the Gospel says that heaven and earth \latin{transíbunt},
``shall pass,'' and the Secret, with the same verb, asks that we pass over,
\latin{transeámus}, to heavenly ``desires.'' Gregory says why the end of the world is no
grief to those who love God: ``Those who love God are bidden to rejoice and
be glad at the world's end, because they soon find him whom they love, while
that which they did not love passes away'' (I.1.3).

The Postcommunion names the sacrament just received. Neither Jerome nor
Augustine mentions the Eucharist at the verse of the eagles: Jerome's
gathering is to the Passion wherever it is read, and Augustine's eagles are
fed ``as it were'' by imitating it. Fromage prays on this Gospel for the elect
``whose youth, like that of the eagle, has been so often renewed by their
receiving thy precious Blood,'' and says that nothing will stay their flight
``once the sacred prey of thy divine Body shall be shown them,'' bringing them
``all together to the banquet of the eternal Pasch'' (vol.~11, pp.~509--510).
In the first interpretation those who ask that ``whatever is evil in our
hearts may be restored'' ask it on the way to the gathering of verse~31. In
the second, what they ask is what Hilary's nursing souls also lack, escape
from their sins. In the third, the prayer has both tenses: the sacrament is
one \latin{quæ súmpsimus}, ``which we have received,'' and the cure is still
asked for, \latin{curétur}, ``may it be cured.''

\subsection{The four senses compared}

The three interpretations read the same Gospel largely with the same Fathers,
and differ in the verse from which each begins and in what each asks of it:
where the end leads, how to live before it, and why it can be trusted. The first weighs
most in the allegorical and anagogical senses. Its allegory makes the body
the crucified and risen Christ in Jerome's, Hilary's and Augustine's three
ways, and the captivity sin forgiven; its anagogy is Hilary's assembly in the
place of the Passion and Jerome's return of the penitent to the Church. The
second weighs most in the moral sense, the rule of the housetop drawn by
Hilary, Jerome and Bede; its anagogy is Jerome's, that the days are shortened
``lest the faith of believers be shaken by the length of the time'' (on
24:22), and Hilary's, that on the day the one who did not go down is found
neither in the cold of sins nor idle from good works. The third weighs most
in the allegorical and moral senses: its allegory is Chrysostom's generation
of believers and Gregory's word that utters unchanging sentences, its moral
sense belief in what is to come from what is fulfilled and prayer with
forgiveness of others; its anagogy is Chrysostom's spiritual summer after the
tempest and Gregory's eternal sun.

Within the first interpretation, Jerome, Hilary, Augustine and Chrysostom place the gathering
differently. The second and third rest on Fathers who disagree about
method: Hilary's step from the letter of the housetop to the soul, which the
second follows, is the step Chrysostom does not take, and the third is built
chiefly on Chrysostom. Neither leaves the other Father out. The second sets
Chrysostom's war first, and the third uses Hilary on the generation that does
not pass and on the words that remain, though not on the fig tree, since
Hilary's branch is Antichrist and Chrysostom's and Gregory's fig is a pledge
of summer.
